\section{Introduction}

The deployment of autonomous multi-agent systems into real-world operational environments represents one of the most consequential transitions in the history of engineered software. Unlike monolithic single-agent systems, multi-agent architectures distribute decision-making authority across interdependent actors, each capable of initiating actions, invoking tools, spawning sub-agents, and modifying shared state. When these agents operate within well-scoped, low-stakes environments, their behavior remains relatively predictable and their failure modes remain contained. The situation changes markedly when agents are deployed in high-stakes domains—financial transactions, critical infrastructure control, medical decision support, or legal reasoning—where the cost of behavioral deviation or loss of executive oversight can be catastrophic.

The need for a mandatory bail-out mechanism in such systems is not immediately obvious to practitioners who have grown accustomed to designing agentic workflows around success paths. The dominant mental model treats agentic systems as inherently goal-directed: the agent receives an objective, it plans, it executes, and it reports. Under this model, failure is an edge case—an anomalous outcome to be caught by exception handlers or human-in-the-loop oversight. However, this model fails to account for the emergent behaviors that arise when multiple agents interact asynchronously, when tool-calling chains grow beyond the designer's intent, or when the agent's internal state diverges from the designer's assumptions due to non-deterministic planning or environment feedback. In our experience building and deploying agentic systems at scale, these are not edge cases. They are the operative condition.

Consider what happens when a multi-agent system lacks a bail-out mechanism. The most benign outcome is that the system enters a deadlock or livelock state: agents continue cycling through redundant actions, consuming resources without progressing toward the stated objective. More severe is the scenario in which one or more agents begin propagating errors through the system faster than the oversight layer can detect and correct them. In the worst cases, agents reach a state from which they cannot be recalled by any agent or human within the system, because no recall mechanism was ever defined. We term this \textit{uncontrolled agentic cascade}—a condition in which the autonomous system's behavior exceeds the \bounds of any defined intervention pathway. The result is not merely a system outage; it is a failure of accountability itself.

The architectural dimension of this problem is what makes it tractable as a research question. Existing approaches to agent safety and alignment focus primarily on the design of the individual agent: its objective function, its training distribution, its constraint specification. While this work is valuable, it addresses only one axis of a fundamentally multi-dimensional problem. The question of what happens at the \bounds of a multi-agent system—where one agent's authority ends and another's begins, where the system's aggregate behavior escapes the \purpose for which it was instantiated, where no human overseer can intervene—is an architectural question, not merely a behavioral one.

This is the core insight of the present work: \accountability must have an architectural fallback. By this we mean that accountability cannot be delegated solely to the agent layer. It must be instantiated as a structural property of the system itself—a mechanism that exists independently of any individual agent's volition or capability. An accountability architecture that depends entirely on agents behaving correctly in order to remain accountable is not an accountability architecture at all. It is a wish. The bail-out mechanism, as we define it in this paper, is precisely such a structural fallback. It is a mandatory, architecture-level intervention that restores the system to a safe state regardless of the current state of any individual agent or agentic subsystem.

We refer to this framework and its associated mechanisms collectively as \bxthree. The \bxthree framework introduces four primitive constructs that govern agentic system behavior at the architectural level. \textbf{\purpose} defines the teleological boundary of the system: what the system is for, expressed in a form that is both machine-readable and human-interpretable. \textbf{\bounds} defines the operational envelope: the hard limits within which the system is permitted to act, including which actions are categorically prohibited, which resources may be accessed, and which states constitute a mandatory bail-out trigger. \textbf{\fact} defines the evidentiary substrate: the provable record of system actions, states, and decisions that can be used to determine whether the system is operating within its \purpose and \bounds. And the fourth construct, which we describe in detail in subsequent sections, is the mechanism by which \purpose, \bounds, and \fact are enforced at runtime—an architectural construct that we term the \bxthree bail-out layer.

The remainder of this paper is organized as follows. Section 2 reviews related work in multi-agent systems safety, agentic alignment, and architectural fallback mechanisms. Section 3 introduces the \bxthree framework in full, defining each of the four primitive constructs and their formal relationships. Section 4 presents the system architecture for the bail-out layer, including the event taxonomy that triggers mandatory intervention, the state machine that governs the transition to safe state, and the proof-of-concept implementation. Section 5 evaluates the framework through a series of experimental scenarios designed to provoke uncontrolled agentic cascade under varying conditions, measuring the effectiveness of mandatory bail-out against baseline systems lacking architectural fallback. Section 6 discusses the limitations of the current approach, including the costs of mandatory intervention in high-throughput systems and the open question of how \bounds should be specified in domains where the operational envelope is inherently ambiguous. Section 7 concludes with a summary of contributions and a forward-looking discussion of the research agenda that follows from this work.
